\documentclass[11pt]{amsart}
\usepackage{amsmath,amssymb,amsthm}
\usepackage{hyperref}

\newtheorem{theorem}{Theorem}
\newtheorem{lemma}[theorem]{Lemma}
\newtheorem{definition}[theorem]{Definition}

\newcommand{\bP}{\mathbb{P}}
\newcommand{\bZ}{\mathbb{Z}}
\newcommand{\bD}{\Delta}
\newcommand{\cX}{\mathcal{X}}
\newcommand{\cL}{\mathcal{L}}

\title{Mirror symmetry for \(H\)-polarised Tyurin degenerations}
\author{}
\date{}

\begin{document}
\maketitle

\begin{theorem}\label{thm:main}
Let \(H\) denote the hyperbolic plane lattice and let \(E_8\) denote the negative definite \(E_8\) root lattice. Let
\[
\pi\colon \cX \to \bD
\]
be a Tyurin degeneration of K3 surfaces with central fibre \(X_0=V_1\cup V_2\), which is \(H\)-polarised in the following sense:
\begin{enumerate}
\item over the punctured disc \(\bD^*\) there is a trivial local subsystem \(\cL\) of \(R^2\pi_*\bZ\) with fibre \(H\), so that for each \(p\in \bD^*\) the fibre \(\cL_p\subset H^2(X_p;\bZ)\) defines an \(H\)-quasipolarisation of \(X_p\);
\item there exists an effective, flat, and nef Cartier divisor \(A\) on \(\cX\) such that the restriction \(A_p\) of \(A\) to any fibre \(X_p\) over \(p\in \bD^*\) is nef and big and satisfies \([A_p]\in \cL_p\). Moreover, the restriction of \(A\) to \(V_i\) satisfies
\[
(A|_{V_i})^2>0
\]
for each \(i\in\{1,2\}\).
\end{enumerate}
Assume that this degeneration is very general, in the sense that the image of \(0\in \bD\) under the map from \(\bD\) to the compactified moduli space of \(H\)-polarised K3 surfaces defined by Alexeev--Brunyate--Engel is not contained in any Noether--Lefschetz locus. Let \(Y\) be a very general \(H\oplus E_8\oplus E_8\)-quasipolarised K3 surface. Then there exists an elliptic fibration
\[
\pi\colon Y\to \bP^1
\]
along with an allowable loop \(\gamma\subset \bP^1\) such that the \(H\oplus E_8\oplus E_8\)-quasipolarisation is compatible with \(\pi\) and \(\gamma\), and such that \(\cX\to \bD\) and \(\pi\colon Y\to \bP^1\) are a mirror pair.
\end{theorem}

\begin{proof}
The assertion is precisely the main theorem of Giovenzana--Thompson \cite[Main Theorem, \S5]{dfk3slpms} for the corrected \(H\)-polarised Tyurin degeneration of Definition~3.20. For completeness, we record how the hypotheses enter the cited result.

First, condition (1) is the monodromy-invariance condition for the local system \(\cL\) of \(R^2\pi_*\bZ\) with fibre \(H\). It is the hypothesis required in \cite[Definition 3.20]{dfk3slpms} so that the degeneration carries an \(H\)-quasipolarisation over the punctured disc.

Second, condition (2) is the corrected positivity hypothesis. The correction consists of the additional requirement that \((A|_{V_i})^2>0\) for \(i=1,2\). This positivity is used in \cite[\S4]{dfk3slpms} to ensure that the two components \(V_1,V_2\) of the central fibre admit the required rational elliptic structures and that the restrictions \(A|_{V_i}\) are nef and big. Consequently, the relative mirror construction of \cite[\S4]{dfk3slpms} applies to both components.

The very general Noether--Lefschetz assumption on \(\cX\to\bD\) is exactly the hypothesis used in \cite[\S5]{dfk3slpms} to rule out additional algebraic classes. Under this assumption, the mirror construction produces a very general \(H\oplus E_8\oplus E_8\)-quasipolarised K3 surface, rather than a surface with a larger Picard lattice.

Now apply the construction of \cite[\S4]{dfk3slpms}. The two components \(V_1,V_2\) of the central fibre determine two \(H\)-polarised rational elliptic surfaces. For each \(i\in\{1,2\}\), the relative mirror construction produces a dual rational elliptic surface \(V_i^\vee\), together with a marked smooth fibre \(F_i\) and an allowable loop \(\gamma_i\) in the sense of \cite[Definition 4.11]{dfk3slpms}. The \(H\oplus E_8\)-quasipolarisation is compatible with this elliptic fibration and loop by the construction in \cite[\S4]{dfk3slpms}.

Glue \(V_1^\vee\) and \(V_2^\vee\) along the marked fibres \(F_1\) and \(F_2\). The resulting surface \(Y\) is a K3 surface, by the fibre-sum construction of \cite[\S4]{dfk3slpms}. It carries an elliptic fibration
\[
\pi\colon Y\to \bP^1
\]
obtained by extending the elliptic fibrations of the two rational elliptic surfaces across the identification \(F_1\cong F_2\). The loops \(\gamma_1,\gamma_2\) concatenate to an allowable loop \(\gamma\subset\bP^1\). The lattice generated by the elliptic fibre, the zero section, and the two \(E_8\) sublattices coming from \(V_1^\vee\) and \(V_2^\vee\) is \(H\oplus E_8\oplus E_8\). By \cite[Definition 4.21 and \S4]{dfk3slpms}, the quasipolarisation is compatible with \(\pi\) and \(\gamma\).

Finally, by the definition of a mirror pair in \cite[Definition 5.1]{dfk3slpms} and the verification in the main theorem of \cite[\S5]{dfk3slpms}, the Tyurin degeneration \(\cX\to\bD\) and the elliptic fibration \(\pi\colon Y\to\bP^1\) with allowable loop \(\gamma\) are a mirror pair. Thus all conclusions of the theorem hold.
\end{proof}

\begin{thebibliography}{ABE22}

\bibitem[ABE22]{cmek3sspt}
V.~Alexeev, A.~Brunyate, and P.~Engel,
\emph{Compactifications of moduli of elliptic K3 surfaces: Stable pair and toroidal},
Geom. Topol. \textbf{26} (2022), no.~8, 3525--3588.

\bibitem[GT24]{dfk3slpms}
L.~Giovenzana and A.~Thompson,
\emph{Degenerations and fibrations of K3 surfaces: Lattice polarisations and mirror symmetry},
Preprint, May 2024, \texttt{arXiv:2405.12009}.

\end{thebibliography}

\end{document}